% GENERATED FILE. Edit canonical lessons, then run npm run generate:books.
% canonical-source-hash: fnv1a64:4eddbbcc1eeeeb5a
% canonical-lessons: PA-C179-samaa, PA-C179-savere, PA-C179-khullha, PA-C179-dupahir, PA-C179-band, PA-C179-status-review, PA-C179-payoff

\chapter{Morning Open, Afternoon Closed}
\label{ch:pa-a2-schedule-status}
\hlchaptermodality{179}

\begin{chapteropening}
\textbf{By the end of this chapter:} \emph{I can read and give one simple public-hours note with morning and afternoon times and open and closed statuses.}

Retrieve a two-sentence public-hours note through listening, speaking, reading, and bounded Gurmukhi writing after two atom-free review steps.
\end{chapteropening}

\section[\texorpdfstring{\pa{ਸਮਾਂ} (samaa)}{ਸਮਾਂ (samaa)}]{\pa{ਸਮਾਂ} — one word for time}

\label{lesson:PA-C179-samaa}



\begin{quote}
Say the familiar schedule once: \textbf{\pa{ਅਮਨ} \pa{ਨਾਲ} \pa{ਅੱਠ} \pa{ਵਜੇ} \pa{ਤੋਂ} \pa{ਦਸ} \pa{ਵਜੇ} \pa{ਤੱਕ} \pa{ਕੰਮ} \pa{ਕਰੋ।} \pa{ਦਸ} \pa{ਵਜੇ} \pa{ਤੋਂ} \pa{ਬਾਅਦ} \pa{ਘਰ} \pa{ਆਉਣਾ।}}
\end{quote}

\begin{grammarlens}[title={Grammar Lens: one schedule noun}]
\begin{quote}
\textbf{\pa{ਸਮਾਂ}} — \emph{samā̃} — time
\end{quote}

This lesson gives one useful schedule word. Keep it small: it does not add a new clock frame.
\end{grammarlens}

\subsection*{Your turn}
\begin{itemize}
\raggedright
  \item \emph{Hear:} \emph{samā̃}; YOU SAY: \textbf{\pa{ਸਮਾਂ}}
  \item \emph{Say it:} \textbf{\pa{ਸਮਾਂ}}
\end{itemize}

\begin{tcolorbox}[breakable,colback=teal!4,colframe=teal!35!black,title={Before you move on}]
Say \emph{time}. (\textbf{\pa{ਸਮਾਂ}}.)

Source: NCERT, \href{https://www.ncertbooks.net/textbook/pdf/bmnjm109.pdf}{Punjabi activity material}, uses \textbf{\pa{ਸਮਾਂ}} in day-part contexts.
\end{tcolorbox}
\section[\texorpdfstring{\pa{ਸਵੇਰੇ} \pa{ਅੱਠ} \pa{ਵਜੇ} (savere atth vaje)}{ਸਵੇਰੇ ਅੱਠ ਵਜੇ (savere atth vaje)}]{\pa{ਸਵੇਰੇ} \pa{ਅੱਠ} \pa{ਵਜੇ} — place the morning before the hour}

\label{lesson:PA-C179-savere}



\begin{quote}
Say \textbf{\pa{ਸਮਾਂ}}. Then retrieve \textbf{\pa{ਦਸ} \pa{ਵਜੇ} \pa{ਤੋਂ} \pa{ਬਾਅਦ}}.
\end{quote}

\begin{grammarlens}[title={Grammar Lens: morning precedes the hour}]
\begin{quote}
\textbf{\pa{ਸਵੇਰੇ} \pa{ਅੱਠ} \pa{ਵਜੇ}} — \emph{savere aṭṭh vaje} — at eight in the morning
\end{quote}

Put \textbf{\pa{ਸਵੇਰੇ}} before the familiar whole-hour phrase. This lesson owns that bounded day-part use.
\end{grammarlens}

\subsection*{Your turn}
\begin{itemize}
\raggedright
  \item \emph{Hear:} \emph{savere aṭṭh vaje}; YOU SAY: \textbf{\pa{ਸਵੇਰੇ} \pa{ਅੱਠ} \pa{ਵਜੇ}}
  \item \emph{Read it:} \textbf{\pa{ਸਵੇਰੇ} \pa{ਅੱਠ} \pa{ਵਜੇ}}
\end{itemize}

\begin{tcolorbox}[breakable,colback=teal!4,colframe=teal!35!black,title={Before you move on}]
Say \emph{at eight in the morning}, then say \emph{time}.

Source: UK Government, \href{https://assets.publishing.service.gov.uk/media/5d947f3340f0b65e64bf87ea/Pack\_B\_-\_Punjabi\_-\_Web.pdf}{Punjabi asylum accommodation guide}, uses \textbf{\pa{ਸਵੇਰੇ}} in public-hours language.
\end{tcolorbox}
\section[\texorpdfstring{\pa{ਬਾਗ਼} \pa{ਸਵੇਰੇ} \pa{ਖੁੱਲ੍ਹਾ} \pa{ਹੈ} (baagh savere …)}{ਬਾਗ਼ ਸਵੇਰੇ ਖੁੱਲ੍ਹਾ ਹੈ (baagh savere …)}]{\pa{ਬਾਗ਼} \pa{ਸਵੇਰੇ} \pa{ਖੁੱਲ੍ਹਾ} \pa{ਹੈ} — one open status}

\label{lesson:PA-C179-khullha}



\begin{quote}
Read \textbf{\pa{ਬਾਗ਼}}. Then say \textbf{\pa{ਸਵੇਰੇ}}.
\end{quote}

\begin{grammarlens}[title={Grammar Lens: one open status}]
\begin{quote}
\textbf{\pa{ਬਾਗ਼} \pa{ਸਵੇਰੇ} \pa{ਖੁੱਲ੍ਹਾ} \pa{ਹੈ}} — \emph{bāġ savere khullhā hai} — the garden is open in the morning
\end{quote}

Here \textbf{\pa{ਖੁੱਲ੍ਹਾ}} gives the open status. Keep the whole model together.
\end{grammarlens}

\subsection*{Your turn}
\begin{itemize}
\raggedright
  \item \emph{Hear:} the open-status model; YOU SAY: \textbf{\pa{ਬਾਗ਼} \pa{ਸਵੇਰੇ} \pa{ਖੁੱਲ੍ਹਾ} \pa{ਹੈ}}
  \item \emph{Read it:} \textbf{\pa{ਬਾਗ਼} \pa{ਸਵੇਰੇ} \pa{ਖੁੱਲ੍ਹਾ} \pa{ਹੈ}}
\end{itemize}

\begin{tcolorbox}[breakable,colback=teal!4,colframe=teal!35!black,title={Before you move on}]
Say the open-status model without looking.

Source: Punjabi University, \href{https://pnarang.weebly.com/uploads/1/4/5/6/14563640/vocab\_f.pdf}{teaching vocabulary}, publishes the \textbf{\pa{ਖੁੱਲ੍ਹਾ}/\pa{ਬੰਦ}} contrast.
\end{tcolorbox}
\section[\texorpdfstring{\pa{ਦੁਪਹਿਰ} \pa{ਦੋ} \pa{ਵਜੇ} (dupahir do vaje)}{ਦੁਪਹਿਰ ਦੋ ਵਜੇ (dupahir do vaje)}]{\pa{ਦੁਪਹਿਰ} \pa{ਦੋ} \pa{ਵਜੇ} — place the afternoon before the hour}

\label{lesson:PA-C179-dupahir}



\begin{quote}
Say \textbf{\pa{ਸਮਾਂ}}. Then retrieve \textbf{\pa{ਬਾਗ਼} \pa{ਸਵੇਰੇ} \pa{ਖੁੱਲ੍ਹਾ} \pa{ਹੈ}}.
\end{quote}

\begin{grammarlens}[title={Grammar Lens: afternoon precedes the hour}]
\begin{quote}
\textbf{\pa{ਦੁਪਹਿਰ} \pa{ਦੋ} \pa{ਵਜੇ}} — \emph{dupahir do vaje} — at two in the afternoon
\end{quote}

Put \textbf{\pa{ਦੁਪਹਿਰ}} before the familiar whole-hour phrase. This lesson owns only this day-part use.
\end{grammarlens}

\subsection*{Your turn}
\begin{itemize}
\raggedright
  \item \emph{Hear:} \emph{dupahir do vaje}; YOU SAY: \textbf{\pa{ਦੁਪਹਿਰ} \pa{ਦੋ} \pa{ਵਜੇ}}
  \item \emph{Read it:} \textbf{\pa{ਦੁਪਹਿਰ} \pa{ਦੋ} \pa{ਵਜੇ}}
\end{itemize}

\begin{tcolorbox}[breakable,colback=teal!4,colframe=teal!35!black,title={Before you move on}]
Say \emph{at two in the afternoon}.

Source: NCERT, \href{https://www.ncertbooks.net/textbook/pdf/bmnjm109.pdf}{Punjabi activity material}, uses \textbf{\pa{ਦੁਪਹਿਰ}} in a day-part context.
\end{tcolorbox}
\section[\texorpdfstring{\pa{ਬਾਗ਼} \pa{ਬੰਦ} \pa{ਹੈ} (baagh band hai)}{ਬਾਗ਼ ਬੰਦ ਹੈ (baagh band hai)}]{\pa{ਬਾਗ਼} \pa{ਬੰਦ} \pa{ਹੈ} — one closed status}

\label{lesson:PA-C179-band}



\begin{quote}
Retrieve \textbf{\pa{ਬਾਗ਼} \pa{ਸਵੇਰੇ} \pa{ਖੁੱਲ੍ਹਾ} \pa{ਹੈ}}.
\end{quote}

\begin{grammarlens}[title={Grammar Lens: one closed status}]
\begin{quote}
\textbf{\pa{ਬਾਗ਼} \pa{ਬੰਦ} \pa{ਹੈ}} — \emph{bāġ band hai} — the garden is closed
\end{quote}

Here \textbf{\pa{ਬੰਦ}} gives the contrasting closed status.
\end{grammarlens}

\subsection*{Your turn}
\begin{itemize}
\raggedright
  \item \emph{Hear:} the closed-status model; YOU SAY: \textbf{\pa{ਬਾਗ਼} \pa{ਬੰਦ} \pa{ਹੈ}}
  \item \emph{Read it:} \textbf{\pa{ਬਾਗ਼} \pa{ਬੰਦ} \pa{ਹੈ}}
\end{itemize}

\begin{tcolorbox}[breakable,colback=teal!4,colframe=teal!35!black,title={Before you move on}]
Say the closed model, then the open model.

Sources: Punjabi University, \href{https://dic.learnpunjabi.org/default.aspx?look=\%E0\%A8\%AC\%E0\%A9\%B0\%E0\%A8\%A6}{RCPLT dictionary}, defines \textbf{\pa{ਬੰਦ}} as closed or shut; Punjabi University \href{https://pnarang.weebly.com/uploads/1/4/5/6/14563640/vocab\_f.pdf}{teaching vocabulary} gives the open/closed pair.
\end{tcolorbox}
\section[\texorpdfstring{\pa{ਸਵੇਰੇ} \pa{ਖੁੱਲ੍ਹਾ}; \pa{ਦੁਪਹਿਰ} \pa{ਬੰਦ} (savere khullha …)}{ਸਵੇਰੇ ਖੁੱਲ੍ਹਾ; ਦੁਪਹਿਰ ਬੰਦ (savere khullha …)}]{\pa{ਸਵੇਰੇ} \pa{ਖੁੱਲ੍ਹਾ}; \pa{ਦੁਪਹਿਰ} \pa{ਬੰਦ} — review the two statuses}

\label{lesson:PA-C179-status-review}



\begin{quote}
Say \textbf{\pa{ਸਮਾਂ}}. Read \textbf{\pa{ਬਾਗ਼} \pa{ਸਵੇਰੇ} \pa{ਖੁੱਲ੍ਹਾ} \pa{ਹੈ}} and \textbf{\pa{ਬਾਗ਼} \pa{ਬੰਦ} \pa{ਹੈ}}. Then say \textbf{\pa{ਦੁਪਹਿਰ} \pa{ਦੋ} \pa{ਵਜੇ}}.
\end{quote}

\subsection*{Your turn}
\begin{itemize}
\raggedright
  \item \emph{Hear:} morning; YOU SAY: \textbf{\pa{ਸਵੇਰੇ}}
  \item \emph{Hear:} afternoon; YOU SAY: \textbf{\pa{ਦੁਪਹਿਰ}}
  \item \emph{Hear:} open; YOU SAY: \textbf{\pa{ਖੁੱਲ੍ਹਾ}}
  \item \emph{Hear:} closed; YOU SAY: \textbf{\pa{ਬੰਦ}}
\end{itemize}

\begin{tcolorbox}[breakable,colback=teal!4,colframe=teal!35!black,title={Before you move on}]
Retrieve all five new forms without looking.

This is reinforcement inside the second prerequisite tranche. It does not unblock Punjabi A2 mock 1 and is not an A2 mock, A2 readiness, or A2 attainment.
\end{tcolorbox}
\section[\texorpdfstring{\pa{ਬਾਗ਼} \pa{ਸਵੇਰੇ} \pa{ਅੱਠ} \pa{ਵਜੇ} … (baagh savere …)}{ਬਾਗ਼ ਸਵੇਰੇ ਅੱਠ ਵਜੇ … (baagh savere …)}]{\pa{ਬਾਗ਼} \pa{ਸਵੇਰੇ} \pa{ਅੱਠ} \pa{ਵਜੇ} \pa{ਖੁੱਲ੍ਹਾ} \pa{ਹੈ।} \pa{ਦੁਪਹਿਰ} \pa{ਦੋ} \pa{ਵਜੇ} \pa{ਬੰਦ} \pa{ਹੈ।}}

\label{lesson:PA-C179-payoff}



\begin{quote}
Retrieve \textbf{\pa{ਸਮਾਂ}}, then the morning, afternoon, open, and closed forms.
\end{quote}

\subsection*{Your turn: four-skill payoff}
\begin{quote}
\textbf{\pa{ਬਾਗ਼} \pa{ਸਵੇਰੇ} \pa{ਅੱਠ} \pa{ਵਜੇ} \pa{ਖੁੱਲ੍ਹਾ} \pa{ਹੈ।} \pa{ਦੁਪਹਿਰ} \pa{ਦੋ} \pa{ਵਜੇ} \pa{ਬੰਦ} \pa{ਹੈ।}}
\end{quote}

\begin{itemize}
\raggedright
  \item \emph{Hear:} the two-sentence hours note; YOU SAY: \textbf{\pa{ਬਾਗ਼} \pa{ਸਵੇਰੇ} \pa{ਅੱਠ} \pa{ਵਜੇ} \pa{ਖੁੱਲ੍ਹਾ} \pa{ਹੈ।} \pa{ਦੁਪਹਿਰ} \pa{ਦੋ} \pa{ਵਜੇ} \pa{ਬੰਦ} \pa{ਹੈ।}}
  \item \emph{Read it:} \textbf{\pa{ਬਾਗ਼} \pa{ਸਵੇਰੇ} \pa{ਅੱਠ} \pa{ਵਜੇ} \pa{ਖੁੱਲ੍ਹਾ} \pa{ਹੈ।} \pa{ਦੁਪਹਿਰ} \pa{ਦੋ} \pa{ਵਜੇ} \pa{ਬੰਦ} \pa{ਹੈ।}}
  \item \emph{Write it:} \textbf{\pa{ਬਾਗ਼} \pa{ਸਵੇਰੇ} \pa{ਅੱਠ} \pa{ਵਜੇ} \pa{ਖੁੱਲ੍ਹਾ} \pa{ਹੈ।} \pa{ਦੁਪਹਿਰ} \pa{ਦੋ} \pa{ਵਜੇ} \pa{ਬੰਦ} \pa{ਹੈ।}} after covering the model
\end{itemize}

\begin{tcolorbox}[breakable,colback=teal!4,colframe=teal!35!black,title={Before you move on}]
Rebuild the hours note without looking, then name the five new forms.

This is the second prerequisite tranche exposed by Punjabi A2 mock drafting. It does not unblock Punjabi A2 mock 1 by itself. It is not an A2 mock, not A2 readiness or attainment, and not external or human validation.

Sources: UK Government, \href{https://assets.publishing.service.gov.uk/media/5d947f3340f0b65e64bf87ea/Pack\_B\_-\_Punjabi\_-\_Web.pdf}{Punjabi asylum accommodation guide}, for public-hours language; Punjabi University, \href{https://pnarang.weebly.com/uploads/1/4/5/6/14563640/vocab\_f.pdf}{teaching vocabulary}, for the open/closed pair.
\end{tcolorbox}
